% Teil II: der durchgerechnete Weg.

\subsection*{Gegeben und gesucht}
Gegeben: geographische Breite $\vartheta=47^\circ$ nördlich,
Fliesgeschwindigkeit $v=\qty{0.5}{m/s}$, $\Omega\approx\qty{7.27e-5}{rad/s}$,
$g=\qty{9.81}{m/s^2}$, Flussbreite $B=\qty{1}{km}$.
Gesucht: (a) drei Unterscheidungsmerkmale; (b) Coriolisbeschleunigung;
(c) betroffenes Ufer und Höhenunterschied des Wassers.

\subsection*{Rechnung}
\textbf{a) Eigenschaften der Trägheitskräfte.}
\begin{itemize}
  \item Sie sind \emph{proportional zur Masse} (wie echte Kräfte) und
       werten sofort in die Bewegungsgleichung ein, sie verursachen aber
       keine Gegenkraft auf einen anderen Körper (kein drittes
       Newtonsches Axiom).
  \item Sie treten \emph{nur in beschleunigten/rotierenden Systemen} auf und
       verschwinden im Inertialsystem identisch.
  \item Sie haben \emph{keine physikalische Quelle}: keine Wechselwirkung
       mit einem Partner; messbar nur relativ zum Bezugssystem
       (Fliehkraft, Corioliskraft, eulersche Kraft).
\end{itemize}

\textbf{b) Coriolisbeschleunigung.} Für eine horizontale Strömung ist nur die
zur Erdfläche senkrechte Projektion von $\vec\Omega$ wirksam:
\begin{equation}
  a_C = 2\,\Omega\,\sin\vartheta\, v
  = 2\cdot 7.27e-5\cdot \sin47^{\circ}\cdot 0.5\ \si{m/s^2}
  \approx \qty{5.3e-5}{m\,s^{-2}}.
\end{equation}

\textbf{c) Ufer und Wasserspiegel.} Der Fluss fliesst auf der Nordhalbkugel
von Norden nach Süden; die Coriolisbeschleunigung zeigt senkrecht zur
Bewegungsrichtung nach \emph{rechts}, also nach Westen. Das Wasser wird an das
\emph{westliche} Ufer gedrückt (dort fliesst es höher, das Ufer wird stärker
ausgewaschen). Höhenunterschied bei $B=\qty{1}{km}$:
\begin{equation}
  \Delta h = B\,\frac{a_C}{g}
  = \frac{1000\cdot 5.3e-5}{9.81}\ \si{m}
  \approx \qty{5.4e-3}{m} \approx \qty{5}{mm}.
\end{equation}

\begin{ergebnis}
(b) $a_C\approx\qty{5.3e-5}{m/s^2}$. (c) Das Wasser hebt sich am
\emph{westlichen} Ufer um $\Delta h
=B\,a_C/g\approx\qty{5}{mm}$ über das östliche; dort wird das Ufer
überproportional abgetragen.
\end{ergebnis}

\begin{kontrolle}
\textbf{Einheiten:} $a_C$: $\mathrm{rad/s}\cdot m/s = m/s^2$; $\Delta h$:
$m\cdot(m/s^2)/(m/s^2)=m$. \checkmark

\textbf{Grössenordnung:} $a_C\approx 5e-5\,m/s^2$ (rund tausendfach kleiner als
$g$); über $1$\,km nur $\approx5$\,mm. Über sehr lange Fliessstrecken summieren
sich die Effekte - genau das sieht man an langen grossen Flüssen.
\end{kontrolle}

\begin{merksatz}
Corioliskraft $=-2m\,\vec\Omega\times\vec v$ lenkt auf der Nordhalbkugel
Bewegungen nach rechts ab (bezogen auf die Bewegungsrichtung). Horizontale
Strömungen spüren wegen der Erdkrümmung nur $\Omega\sin\vartheta$.
\end{merksatz}